\section{Soluções Existentes}

\vspace{-0.4cm}

\noindent\rule{\textwidth}{0.7pt}

<Se optar pelo capítulo de Soluções Existentes: Identificar e descrever soluções existentes para o problema de pesquisa definido. Detalhar como estas soluções atendem o problema, em que aspectos ou necessidades elas atendem. E descrever em que as soluções existentes deixam a desejar, ou seja, no que não atendem ao problema. Ao final realizar uma análise comparativa entre as soluções existentes. Mínimo de 5 soluções>

<É obrigatório organizar as soluções existentes em uma estrutura de tópicos.>

\subsection{Solução Existente 1}

<Em um texto corrido, faça a apresentação detalhada da solução. Explique o objetivo da solução, como usar, a aplicabilidade, histórico de sucesso, quando foi criado, quem criou, e o que mais julgar se necessário para explicar a solução.>

\textcolor{red}{
<Os parágrafos produzidos precisam fazer uso de citações.
O recomendável é 1 citação (direta ou indireta)
para cada 2 parágrafos produzidos.>
}

\subsubsection{Aspectos Positivos de Atendimento ao Problema}
<Texto do Tópico> É para detalhar e explicar cada aspecto e não somente citar eles.

\subsubsection{Aspectos Negativos de Atendimento ao Problema}
<Texto do Tópico> É para detalhar e explicar cada aspecto e não somente citar eles.

\subsection{Solução Existente 2}

<Em um texto corrido, faça a apresentação detalhada da solução. Explique o objetivo da solução, como usar, a aplicabilidade, histórico de sucesso, quando foi criado, quem criou, e o que mais julgar se necessário para explicar a solução.>

\textcolor{red}{
<Os parágrafos produzidos precisam fazer uso de citações.
O recomendável é 1 citação (direta ou indireta)
para cada 2 parágrafos produzidos.>
}

\subsubsection{Aspectos Positivos de Atendimento ao Problema}
<Texto do Tópico> É para detalhar e explicar cada aspecto e não somente citar eles.

\subsubsection{Aspectos Negativos de Atendimento ao Problema}
<Texto do Tópico> É para detalhar e explicar cada aspecto e não somente citar eles.

\subsection{Solução Existente 3}

<Em um texto corrido, faça a apresentação detalhada da solução. Explique o objetivo da solução, como usar, a aplicabilidade, histórico de sucesso, quando foi criado, quem criou, e o que mais julgar se necessário para explicar a solução.>

\textcolor{red}{
<Os parágrafos produzidos precisam fazer uso de citações.
O recomendável é 1 citação (direta ou indireta)
para cada 2 parágrafos produzidos.>
}

\subsubsection{Aspectos Positivos de Atendimento ao Problema}
<Texto do Tópico> É para detalhar e explicar cada aspecto e não somente citar eles.

\subsubsection{Aspectos Negativos de Atendimento ao Problema}
<Texto do Tópico> É para detalhar e explicar cada aspecto e não somente citar eles.

\subsection{Solução Existente 4}

<Em um texto corrido, faça a apresentação detalhada da solução. Explique o objetivo da solução, como usar, a aplicabilidade, histórico de sucesso, quando foi criado, quem criou, e o que mais julgar se necessário para explicar a solução.>

\textcolor{red}{
<Os parágrafos produzidos precisam fazer uso de citações.
O recomendável é 1 citação (direta ou indireta)
para cada 2 parágrafos produzidos.>
}

\subsubsection{Aspectos Positivos de Atendimento ao Problema}
<Texto do Tópico> É para detalhar e explicar cada aspecto e não somente citar eles.

\subsubsection{Aspectos Negativos de Atendimento ao Problema}
<Texto do Tópico> É para detalhar e explicar cada aspecto e não somente citar eles.

\subsection{Solução Existente 5}

<Em um texto corrido, faça a apresentação detalhada da solução. Explique o objetivo da solução, como usar, a aplicabilidade, histórico de sucesso, quando foi criado, quem criou, e o que mais julgar se necessário para explicar a solução.>

\textcolor{red}{
<Os parágrafos produzidos precisam fazer uso de citações.
O recomendável é 1 citação (direta ou indireta)
para cada 2 parágrafos produzidos.>
}

\subsubsection{Aspectos Positivos de Atendimento ao Problema}
<Texto do Tópico> É para detalhar e explicar cada aspecto e não somente citar eles.

\subsubsection{Aspectos Negativos de Atendimento ao Problema}
<Texto do Tópico> É para detalhar e explicar cada aspecto e não somente citar eles.

\subsection{ Análise Comparativa das Soluções Existentes}
<Analise as 5 soluções existentes comparativamente sob a óticas dos critérios pertinentes aos problemas de pesquisa. Para isto, crie uma tabela para apresentar os pontos em comum e divergentes entre as soluções selecionadas.>


\begin{table}[htbp]
\centering
\caption{Tabela comparativa das soluções existentes}
\label{tab:comparativo-solucoes}
\renewcommand{\arraystretch}{1.2}
\begin{tabular}{|p{3.2cm}|p{2.2cm}|p{2.2cm}|p{2.2cm}|p{2.2cm}|p{2.2cm}|}
\hline
\textbf{Critério} & \textbf{Solução 1} & \textbf{Solução 2} & \textbf{Solução 3} & \textbf{Solução 4} & \textbf{Solução 5} \\
\hline
\textbf{Critério 1} & \textit{(preencher)} & \textit{(preencher)} & \textit{(preencher)} & \textit{(preencher)} & \textit{(preencher)} \\
\hline
\textbf{Critério 2} & \textit{(preencher)} & \textit{(preencher)} & \textit{(preencher)} & \textit{(preencher)} & \textit{(preencher)} \\
\hline
\textbf{Critério 3} & \textit{(preencher)} & \textit{(preencher)} & \textit{(preencher)} & \textit{(preencher)} & \textit{(preencher)} \\
\hline
\textbf{Critério 4} & \textit{(preencher)} & \textit{(preencher)} & \textit{(preencher)} & \textit{(preencher)} & \textit{(preencher)} \\
\hline
\textbf{Critério N} & \textit{(preencher)} & \textit{(preencher)} & \textit{(preencher)} & \textit{(preencher)} & \textit{(preencher)} \\
\hline
\end{tabular}
\end{table}

<Após a construção da tabela de soluções e critérios é feita a análise interpretativa.>

<Você precisa explicar [e não descrever] a tabela. Pense nessas perguntas: Uma solução não possuir um determinado conjunto de critério significa o que? Um critério não ser atendido significa o que?>